\emph{অপেক্ষক}~$f\colon A \to B$ তার \emph{সংজ্ঞাক্ষেত্র}~$A$-এর
প্রত্যেক উপাদানকে \emph{সহসংজ্ঞাক্ষেত্র}~$B$-এর একটি নির্দিষ্ট
উপাদানে চিত্রিত করে। $x \in A$ হলে, যে $y$-তে~$f$, $x$-কে
চিত্রিত করে, তাকে \emph{আর্গুমেন্ট}~$x$-এর জন্য~$f$-এর
\emph{মান}~$f(x)$ বলি। $A$ যুগলের সেট হলে অপেক্ষক~$f$-কে
দুটি আর্গুমেন্ট গ্রহণকারী অপেক্ষক হিসেবে ভাবতে পারি। $f$-এর
\emph{বিস্তৃতি}~$\ran{f}$ হলো~$B$-এর সেই উপসেট যা~$f$-এর সব
মান নিয়ে গঠিত।

$\ran{f} = B$ হলে $f$-কে \emph{সমাপতিত} বলে। মান~$f(x)$ অনন্য
এই অর্থে যে $f$, $x$-কে কেবল একটি $f(x)$-এ চিত্রিত করে, কখনোই
একাধিকটিতে নয়। $f(x)$ যদি এই অর্থেও অনন্য হয় যে দুটি ভিন্ন
আর্গুমেন্ট একই মানে চিত্রিত হয় না, তবে $f$-কে \emph{একৈক} বলে।
একৈক এবং সমাপতিত দুই-ই এমন অপেক্ষককে \emph{একৈক সমাপতিত} বলে।

একৈক সমাপতিত অপেক্ষকের অনন্য \emph{বিপরীত অপেক্ষক}~$f^{-1}$ থাকে।
অপেক্ষকগুলিকেও পরপর জোড়া যায়। অপেক্ষক $(g \circ f)$ হলো $f$-এর
সঙ্গে $g$-এর \emph{মিশ্রণ}। একৈক অপেক্ষকের মিশ্রণ একৈক, আর
সমাপতিত অপেক্ষকের মিশ্রণ সমাপতিত। $(f^{-1} \circ f)$ হলো অভেদ
অপেক্ষক।

প্রত্যেক $x \in A$-এর জন্য $f$-এর একটি মান থাকতে হবে—এই শর্ত
শিথিল করলে \emph{আংশিক অপেক্ষকের} ধারণা পাওয়া যায়। $f\colon
A \pto B$ আংশিক হলে, আর্গুমেন্ট~$x$-এর জন্য $f$-এর মান থাকলে
বলি $f(x)$ \emph{সংজ্ঞায়িত}, $f(x) \fdefined$। অন্যথায় বলি
$f(x)$ \emph{অসংজ্ঞায়িত}, $f(x) \fundefined$। যেকোনো (আংশিক)
অপেক্ষক~$f$-এর সঙ্গে $f$-এর \emph{গ্রাফ}~$R_f$ যুক্ত থাকে। এই
সম্পর্কটি থাকে ঠিক তখনই, যখন $f(x) = y$।
